\documentclass[12pt, a4paper]{article}

% Pakiety
\usepackage[utf8]{inputenc}
\usepackage[T1]{fontenc}
\usepackage[polish]{babel}
\usepackage{amsmath, amssymb}
\usepackage{graphicx}
\usepackage{float}
\usepackage{geometry}
\usepackage{hyperref}
\usepackage{booktabs}
\usepackage{caption}
\usepackage{listings}
\usepackage{xcolor}
\usepackage{parskip}

\geometry{
  a4paper,
  left=2.5cm,
  right=2.5cm,
  top=2.5cm,
  bottom=2.5cm
}

% Zakaz dzielenia wyrazów
\hyphenpenalty=10000
\exhyphenpenalty=10000
\emergencystretch=3em

\hypersetup{
  colorlinks=true,
  linkcolor=blue,
  urlcolor=blue,
  citecolor=blue
}

% Zmiana nazwy float figure: Rysunek -> Wykres
\addto\captionspolish{\renewcommand{\figurename}{Wykres}}

% Styl listingów
\lstset{
  language=Python,
  basicstyle=\ttfamily\small,
  keywordstyle=\color{blue}\bfseries,
  commentstyle=\color{gray}\itshape,
  stringstyle=\color{orange!80!black},
  numbers=left,
  numberstyle=\tiny\color{gray},
  numbersep=8pt,
  frame=single,
  breaklines=true,
  breakatwhitespace=true,
  captionpos=b,
  showstringspaces=false,
  columns=flexible,
  keepspaces=true,
  inputencoding=utf8,
  extendedchars=true,
  literate={ą}{{\k{a}}}1 {ę}{{\k{e}}}1 {ó}{{\'{o}}}1
            {ś}{{\'{s}}}1 {ł}{{\l{}}}1 {ź}{{\'{z}}}1
            {ż}{{\.z}}1 {ć}{{\'{c}}}1 {ń}{{\'{n}}}1
            {Ą}{{\k{A}}}1 {Ę}{{\k{E}}}1 {Ó}{{\'{O}}}1
            {Ś}{{\'{S}}}1 {Ł}{{L}}1,
}

% Strona tytułowa
\title{
  \textbf{Metody Obliczeniowe w Nauce i Technice} \\[0.5em]
  \large Laboratorium 8 \textendash{} Rozwiązywanie równań nieliniowych
}
\author{Mateusz Stoch}
\date{6 maja 2026}

\begin{document}

\maketitle

% ============================================================
\section{Zadanie 1 \textendash{} Schematy iteracyjne}
% ============================================================

\subsection{Opis zadania}

Dane jest równanie kwadratowe:
\[
  f(x) = x^2 - 3x + 2 = 0
\]
którego pierwiastki to $\alpha_1 = 1$ i $\alpha_2 = 2$. Badamy zbieżność do pierwiastka $\alpha = 2$ czterech równoważnych schematów iteracyjnych:

\begin{align*}
  \varphi_1(x) &= \frac{x^2 + 2}{3} \\
  \varphi_2(x) &= \sqrt{3x - 2} \\
  \varphi_3(x) &= 3 - \frac{2}{x} \\
  \varphi_4(x) &= \frac{x^2 - 2}{2x - 3}
\end{align*}

\subsection{Analiza zbieżności}

Warunkiem zbieżności iteracji $x_{k+1} = \varphi(x_k)$ do punktu stałego $\alpha$ jest $|\varphi'(\alpha)| < 1$. Obliczamy pochodne i sprawdzamy ich wartości w punkcie $\alpha = 2$:

\begin{center}
\begin{tabular}{l|c|c}
\toprule
Schemat & $|\varphi'(2)|$ & Zbieżny \\
\midrule
$\varphi_1 = (x^2+2)/3$ & $1{,}333333$ & NIE \\
$\varphi_2 = \sqrt{3x-2}$ & $0{,}750000$ & TAK \\
$\varphi_3 = 3 - 2/x$ & $0{,}500000$ & TAK \\
$\varphi_4 = (x^2-2)/(2x-3)$ & $0{,}000000$ & TAK \\
\bottomrule
\end{tabular}
\end{center}

\subsection{Wyniki eksperymentalne}

Każdy schemat wykonano przez 10 iteracji, startując z $x_0 = 2{,}5$. Rząd zbieżności wyznaczono empirycznie ze wzoru:
\[
  r = \frac{\ln(\varepsilon_k / \varepsilon_{k+1})}{\ln(\varepsilon_{k-1} / \varepsilon_k)}, \quad \varepsilon_k = |x_k - \alpha|
\]

\begin{lstlisting}[caption={Implementacja schematów iteracyjnych}]
alpha = np.float64(2.0)
phi1 = lambda x: (x**2 + 2) / 3
phi2 = lambda x: np.sqrt(3*x - 2) if 3*x - 2 >= 0 else float('nan')
phi3 = lambda x: 3 - 2/x
phi4 = lambda x: (x**2 - 2) / (2*x - 3)
\end{lstlisting}

Poniżej wybrane iteracje dla każdego schematu:

\begin{center}
\begin{tabular}{l|r|r}
\toprule
Schemat & $x_{10}$ & $\varepsilon_{10}$ \\
\midrule
$\varphi_1$ (rozbieżny) & \texttt{nan} & \texttt{nan} \\
$\varphi_2$ & $2{,}0210522$ & $2{,}11 \times 10^{-2}$ \\
$\varphi_3$ & $2{,}0003256$ & $3{,}26 \times 10^{-4}$ \\
$\varphi_4$ & $2{,}0000000$ & $0{,}00 \times 10^{0}$ \\
\bottomrule
\end{tabular}
\end{center}

Wyznaczone empiryczne rzędy zbieżności:
\begin{itemize}
  \item $\varphi_2$: $r \approx 0{,}992$ (zbieżność liniowa)
  \item $\varphi_3$: $r \approx 0{,}999$ (zbieżność liniowa)
  \item $\varphi_4$: $r \approx 1{,}856$ (zbieżność bliżej kwadratowej)
\end{itemize}

Uwaga: $\varphi_1$ jest rozbieżny, więc wyznaczenie rzędu zbieżności nie ma sensu.

\subsection{Wykresy błędów}

\begin{figure}[H]
  \centering
  \includegraphics[width=0.95\textwidth]{wykres1.png}
  \caption{Błąd względny $|\varepsilon_k|$ w skali logarytmicznej. Lewy panel: wszystkie metody (widać rozbieżność $\varphi_1$). Prawy panel: tylko metody zbieżne.}
  \label{fig:bledy_z1}
\end{figure}

\subsection{Wnioski}

Schemat $\varphi_1$ rozbieżny, bo $|\varphi_1'(2)| = 4/3 > 1$. Schematy $\varphi_2$ i $\varphi_3$ zbiegają liniowo, ponieważ ich pochodne w punkcie stałym są w przedziale $(0, 1)$. Schemat $\varphi_4$ zbiega szybko i odpowiada metodzie Newtona zastosowanej do $f(x)$, stąd jego zbieżność kwadratowa.

% ============================================================
\section{Zadanie 2 \textendash{} Metoda Newtona \textendash{} przypadki problematyczne}
% ============================================================

\subsection{Opis zadania}

Metoda Newtona dla $f(x)$ wyznacza kolejne przybliżenia ze wzoru:
\[
  x_{k+1} = x_k - \frac{f(x_k)}{f'(x_k)}
\]

Sprawdzamy cztery funkcje, dla których metoda Newtona zawodzi przy podanych punktach startowych. Następnie szukamy skutecznych rozwiązań.

\subsection{Przypadek (a): $f(x) = x^3 - 5x$, $x_0 = 1$}

\textbf{Dlaczego zawodzi:} W punkcie $x_0 = 1$ mamy $f'(1) = 3 \cdot 1 - 5 = -2$. Newton wyznacza $x_1 = 1 - (-4)/(-2) = 1 - 2 = -1$, następnie $x_2 = 1$, i tak dalej. Algorytm wpada w cykl i nie zbiega.

\textbf{Rozwiązanie:} Zmiana punktu startowego na $x_0 = 2{,}5$ daje pierwiastek $\sqrt{5} \approx 2{,}23606798$.

\subsection{Przypadek (b): $f(x) = x^3 - 3x + 1$, $x_0 = 1$}

\textbf{Dlaczego zawodzi:} Pochodna $f'(1) = 3 \cdot 1 - 3 = 0$. Newton próbuje podzielić przez zero, co natychmiast kończy iteracje błędem.

\textbf{Rozwiązanie:} Zmiana punktu startowego na $x_0 = 1{,}5$ daje pierwiastek $\approx 1{,}53208889$.

\subsection{Przypadek (c): $f(x) = 2 - x^5$, $x_0 = 0{,}01$}

\textbf{Dlaczego zawodzi:} Pochodna $f'(0{,}01) = -5 \cdot (0{,}01)^4$ jest bardzo mała, co powoduje skok Newtona w odległe miejsce. Iteracja nie wraca w okolice pierwiastka.

\textbf{Rozwiązanie:} Zmiana punktu startowego na $x_0 = 1{,}0$ daje pierwiastek $2^{1/5} \approx 1{,}14869835$.

\subsection{Przypadek (d): $f(x) = x^4 - 4{,}29x^2 - 5{,}29$, $x_0 = 0{,}8$}

\textbf{Dlaczego zawodzi:} Z punktu $x_0 = 0{,}8$ metoda Newtona oscyluje między wartościami, nie zbiegając do żadnego pierwiastka.

\textbf{Rozwiązanie:} Zmiana punktu startowego na $x_0 = 2{,}3$ daje pierwiastek $\approx 2{,}30000000$.

\begin{lstlisting}[caption={Przykładowe naprawione wywołanie dla przypadku (a)}]
f_a  = lambda x: x**3 - 5*x
fp_a = lambda x: 3*x**2 - 5
root_a = optimize.newton(f_a, 2.5, fprime=fp_a)
# wynik: sqrt(5) = 2.23606798
\end{lstlisting}

\subsection{Wnioski}

Metoda Newtona może zawodzić z kilku powodów: zerowa pochodna w punkcie startowym, wpadanie w cykl, lub zbyt mała pochodna powodująca duże skoki. We wszystkich przypadkach wystarczyła zmiana punktu startowego lub, alternatywnie, użycie metody \texttt{scipy.optimize.brentq} jako odporniejszej na takie sytuacje.

% ============================================================
\section{Zadanie 3 \textendash{} Schematy Newtona bez dzielenia}
% ============================================================

\subsection{Część (a): Obliczanie $x = 1/c$}

\textbf{Wyprowadzenie:} Szukamy $x = 1/c$, czyli pierwiastka funkcji $f(x) = \frac{1}{x} - c$. Jej pochodna to $f'(x) = -\frac{1}{x^2}$. Krok Newtona daje:
\[
  x_{k+1} = x_k - \frac{f(x_k)}{f'(x_k)} = x_k - \frac{\frac{1}{x_k} - c}{-\frac{1}{x_k^2}} = x_k + x_k^2\left(\frac{1}{x_k} - c\right) = 2x_k - cx_k^2
\]

\textbf{Schemat iteracji:}
\[
  x_{k+1} = x_k(2 - cx_k) \quad \text{[tylko dodawanie i mnożenie]}
\]

\textbf{Wyniki dla $c = 3$, $x_0 = 0{,}3$:}

\begin{center}
\begin{tabular}{r|r|r}
\toprule
$k$ & $x_k$ & $\varepsilon_k$ \\
\midrule
1 & $0{,}3300000000$ & $3{,}33 \times 10^{-3}$ \\
2 & $0{,}3333000000$ & $3{,}33 \times 10^{-5}$ \\
3 & $0{,}3333333300$ & $3{,}33 \times 10^{-9}$ \\
4 & $0{,}3333333333$ & $5{,}55 \times 10^{-17}$ \\
5+ & $0{,}3333333333$ & $0{,}00 \times 10^{0}$ \\
\bottomrule
\end{tabular}
\end{center}

Po 4 iteracjach osiągamy pełną precyzję \texttt{float64}.

\subsection{Część (b): Obliczanie $x = 1/\sqrt{c}$}

\textbf{Wyprowadzenie:} Szukamy $x = 1/\sqrt{c}$, czyli pierwiastka $f(x) = \frac{1}{x^2} - c$. Pochodna: $f'(x) = -\frac{2}{x^3}$. Krok Newtona:
\[
  x_{k+1} = x_k - \frac{\frac{1}{x_k^2} - c}{-\frac{2}{x_k^3}} = x_k + \frac{x_k}{2}\left(1 - cx_k^2\right) = \frac{x_k(3 - cx_k^2)}{2}
\]

\textbf{Schemat iteracji:}
\[
  x_{k+1} = \frac{x_k(3 - cx_k^2)}{2} \quad \text{[dodawanie, mnożenie, dzielenie przez 2]}
\]

\textbf{Wyniki dla $c = 4$, $x_0 = 0{,}45$:}

\begin{center}
\begin{tabular}{r|r|r}
\toprule
$k$ & $x_k$ & $\varepsilon_k$ \\
\midrule
1 & $0{,}4927500000$ & $7{,}25 \times 10^{-3}$ \\
2 & $0{,}4998430747$ & $1{,}57 \times 10^{-4}$ \\
3 & $0{,}4999999261$ & $7{,}39 \times 10^{-8}$ \\
4 & $0{,}5000000000$ & $1{,}64 \times 10^{-14}$ \\
5+ & $0{,}5000000000$ & $0{,}00 \times 10^{0}$ \\
\bottomrule
\end{tabular}
\end{center}

\subsection{Część (c): Schemat Herona dla $x = \sqrt{c}$}

\textbf{Wyprowadzenie:} Szukamy $x = \sqrt{c}$, czyli pierwiastka $f(x) = x^2 - c$. Pochodna: $f'(x) = 2x$. Krok Newtona:
\[
  x_{k+1} = x_k - \frac{x_k^2 - c}{2x_k} = \frac{x_k + c/x_k}{2} \quad \text{[schemat Herona]}
\]

\textbf{Wada:} schemat wymaga dzielenia przez $x_k$, a nie tylko przez 2. Jeśli $x_k \approx 0$, mamy problem numeryczny, a poza tym instrukcja ogólnego dzielenia jest kosztowna na wielu procesorach.

\textbf{Wyniki dla $c = 2$, $x_0 = 1{,}5$:}

\begin{center}
\begin{tabular}{r|r|r}
\toprule
$k$ & $x_k$ & $\varepsilon_k$ \\
\midrule
1 & $1{,}4166666667$ & $2{,}45 \times 10^{-3}$ \\
2 & $1{,}4142156863$ & $2{,}12 \times 10^{-6}$ \\
3 & $1{,}4142135624$ & $1{,}59 \times 10^{-12}$ \\
4+ & $1{,}4142135624$ & $\leq 2{,}22 \times 10^{-16}$ \\
\bottomrule
\end{tabular}
\end{center}

\subsection{Część (d): Alternatywa dla Herona}

Zamiast bezpośrednio obliczać $\sqrt{c}$ schematem Herona, używamy schematu z części (b) do wyznaczenia $y = 1/\sqrt{c}$, a następnie $\sqrt{c} = c \cdot y$. Całe obliczenie wymaga tylko dodawania, mnożenia i dzielenia przez 2.

\textbf{Schemat:}
\[
  y_{k+1} = \frac{y_k(3 - cy_k^2)}{2}, \quad \text{następnie } \sqrt{c} = c \cdot y
\]

\textbf{Wyniki dla $c = 2$, $y_0 = 0{,}68$:}

\begin{center}
\begin{tabular}{r|r|r|r}
\toprule
$k$ & $1/\sqrt{c}$ & $\sqrt{c}$ & $\varepsilon_k$ \\
\midrule
1 & $0{,}7055680000$ & $1{,}4111360000$ & $3{,}08 \times 10^{-3}$ \\
2 & $0{,}7071017619$ & $1{,}4142035237$ & $1{,}00 \times 10^{-5}$ \\
3 & $0{,}7071067811$ & $1{,}4142135623$ & $1{,}07 \times 10^{-10}$ \\
4+ & $0{,}7071067812$ & $1{,}4142135624$ & $\leq 2{,}22 \times 10^{-16}$ \\
\bottomrule
\end{tabular}
\end{center}

\subsection{Wnioski}

Wszystkie schematy iteracyjne wywodzą się bezpośrednio z metody Newtona. Kluczową obserwacją jest to, że dobierając odpowiednią funkcję $f(x)$ można tak skonstruować schemat, by unikał on kosztownych operacji arytmetycznych. Schemat Herona, choć elegancki, wymaga ogólnego dzielenia, natomiast alternatywa z podpunktu (d) unika tego problemu, zachowując tę samą szybkość zbieżności (kwadratową).

% ============================================================
\section{Zadanie 4 \textendash{} Układ równań nieliniowych metodą Newtona}
% ============================================================

\subsection{Opis zadania}

Rozważamy układ równań nieliniowych:
\[
  \begin{cases}
    x_1^2 + x_2^2 = 1 \\
    x_1^2 - x_2 = 0
  \end{cases}
\]

Dokładne rozwiązanie to:
\[
  x_1 = \sqrt{\frac{\sqrt{5}}{2} - \frac{1}{2}} \approx 0{,}7861513778, \qquad
  x_2 = \frac{\sqrt{5}}{2} - \frac{1}{2} \approx 0{,}6180339887
\]

\subsection{Schemat iteracji Newtona}

Dla układu $\mathbf{F}(\mathbf{x}) = \mathbf{0}$ każdy krok Newtona wyznaczamy rozwiązując układ liniowy:
\[
  J(\mathbf{x}_k)\, \mathbf{s}_{k} = -\mathbf{F}(\mathbf{x}_k), \qquad \mathbf{x}_{k+1} = \mathbf{x}_k + \mathbf{s}_{k}
\]
gdzie $J$ to jakobian układu:
\[
  J = \begin{pmatrix} 2x_1 & 2x_2 \\ 2x_1 & -1 \end{pmatrix}
\]

\begin{lstlisting}[caption={Implementacja metody Newtona dla układu równań}]
def F(xy):
    x1, x2 = xy[0], xy[1]
    return np.array([x1**2 + x2**2 - 1.0, x1**2 - x2])

def J(xy):
    x1, x2 = xy[0], xy[1]
    return np.array([[2.0*x1, 2.0*x2], [2.0*x1, -1.0]])

xy = np.array([1.0, 0.5])
for k in range(10):
    delta = np.linalg.solve(J(xy), -F(xy))
    xy = xy + delta
\end{lstlisting}

\subsection{Wyniki}

Iteracje startując z $\mathbf{x}_0 = (1{,}0,\; 0{,}5)$:

\begin{center}
\begin{tabular}{r|r|r|r}
\toprule
$k$ & $x_1$ & $x_2$ & $\|\mathbf{x}_k - \mathbf{x}^*\|$ \\
\midrule
0 & $1{,}0000000000$ & $0{,}5000000000$ & \textendash{} \\
1 & $0{,}8125000000$ & $0{,}6250000000$ & $2{,}73 \times 10^{-2}$ \\
2 & $0{,}7865918803$ & $0{,}6180555556$ & $4{,}41 \times 10^{-4}$ \\
3 & $0{,}7861515012$ & $0{,}6180339890$ & $1{,}23 \times 10^{-7}$ \\
4 & $0{,}7861513778$ & $0{,}6180339887$ & $9{,}66 \times 10^{-15}$ \\
5+ & $0{,}7861513778$ & $0{,}6180339887$ & $\leq 1{,}11 \times 10^{-16}$ \\
\bottomrule
\end{tabular}
\end{center}

Błąd względny wyznaczonego rozwiązania:
\begin{itemize}
  \item $x_1$: błąd względny $= 0{,}00 \times 10^{0}$
  \item $x_2$: błąd względny $= 1{,}80 \times 10^{-16}$
\end{itemize}

\subsection{Wnioski}

Metoda Newtona dla układów równań zbiega tak samo szybko jak w przypadku skalarnym, czyli kwadratowo. Już po 4 iteracjach błąd spada poniżej progu precyzji maszynowej. Kluczowe jest oczywiście to, że jakobian musi być nieosobliwy w każdym kroku iteracji, a punkt startowy powinien leżeć wystarczająco blisko rozwiązania.

% ============================================================
\section{Podsumowanie}
% ============================================================

Laboratorium 8 dotyczyło numerycznych metod rozwiązywania równań nieliniowych. Najważniejsze wnioski:

\begin{itemize}
  \item Zbieżność iteracji punktu stałego zależy bezpośrednio od wartości $|\varphi'(\alpha)|$: jeśli jest mniejsza od 1, metoda zbiega, w przeciwnym razie rozbieżnie. Schemat odpowiadający metodzie Newtona ($\varphi'(\alpha) = 0$) zbiega kwadratowo.
  \item Metoda Newtona może zawodzić przy nieudanym wyborze punktu startowego (zerowa pochodna, cykl, duży skok). Odpowiedni dobór $x_0$ lub przejście na odporniejszą metodę (np. \texttt{brentq}) rozwiązuje problem.
  \item Sprytny dobór funkcji $f(x)$ w metodzie Newtona pozwala wyeliminować kosztowne operacje arytmetyczne z schematu iteracyjnego, zachowując szybką zbieżność kwadratową.
  \item Rozszerzenie metody Newtona na układy równań wymaga obliczania i odwracania jakobianu, ale zachowuje zbieżność kwadratową.
\end{itemize}

\end{document}
